\documentclass[12pt,a4paper]{article}

% ── Packages ────────────────────────────────────────────────────────────────
\usepackage[T1]{fontenc}
\usepackage[utf8]{inputenc}
\usepackage{lmodern}
\usepackage{microtype}
\usepackage{amsmath,amssymb,amsthm}
\usepackage{mathtools}
\usepackage{enumitem}
\usepackage{geometry}
\usepackage{hyperref}
\usepackage{booktabs}
\usepackage{array}
\usepackage{xcolor}
\usepackage{listings}
\usepackage{fancyhdr}

\geometry{margin=2.5cm}
\setlength{\headheight}{13.60001pt}
\addtolength{\topmargin}{-1.60001pt}

% ── Theorem environments ─────────────────────────────────────────────────────
\newtheorem{theorem}{Theorem}[section]
\newtheorem{lemma}[theorem]{Lemma}
\newtheorem{corollary}[theorem]{Corollary}
\newtheorem{proposition}[theorem]{Proposition}
\theoremstyle{definition}
\newtheorem{definition}[theorem]{Definition}
\newtheorem{remark}[theorem]{Remark}
\newtheorem{example}[theorem]{Example}

% ── Code listings ────────────────────────────────────────────────────────────
\definecolor{codegray}{rgb}{0.96,0.96,0.96}
\definecolor{codeblue}{rgb}{0.1,0.2,0.6}
\definecolor{codegreencomment}{rgb}{0.13,0.55,0.13}

\lstset{
  backgroundcolor=\color{codegray},
  basicstyle=\ttfamily\small,
  keywordstyle=\color{codeblue}\bfseries,
  commentstyle=\color{codegreencomment}\itshape,
  breaklines=true,
  frame=single,
  framerule=0pt,
  rulecolor=\color{gray!30},
  xleftmargin=6pt, xrightmargin=6pt,
}

\lstdefinelanguage{Lean4}{
  keywords={theorem,lemma,def,axiom,by,have,obtain,exact,intro,rw,push_cast,
            ring,decide,linarith,nlinarith,rcases,calc,simp,fin_cases,revert,
            import,where,fun,if,then,else,let,in,match,with,show,apply,refine,
            constructor,left,right,norm_num,norm_cast,exact_mod_cast,
            unfold,push_neg,omega,Set,Finset,open,use,private,rintro,mul_left_cancel₀},
  sensitive=true,
  morecomment=[l]{--},
  morecomment=[s]{/-}{-/},
  morestring=[b]",
  keywordstyle=\color{codeblue}\bfseries,
  commentstyle=\color{codegreencomment}\itshape,
}

% ── Macros ───────────────────────────────────────────────────────────────────
\newcommand{\Z}{\mathbb{Z}}
\newcommand{\Q}{\mathbb{Q}}
\newcommand{\N}{\mathbb{N}}

% ── Header / Footer ──────────────────────────────────────────────────────────
\pagestyle{fancy}
\fancyhf{}
\rhead{\small\thepage}
\lhead{\small Infinitely Many Solutions to $x^5 + y^4 = 3z^2$}
\renewcommand{\headrulewidth}{0.4pt}

% ── Title ────────────────────────────────────────────────────────────────────
\title{%
  \Large\bfseries Infinitely Many Integer Solutions to\\[6pt]
  \Huge $x^5 + y^4 = 3z^2$\\[10pt]
  \large Complete Analysis with Lean~4 Formalisation
}
\author{%
  \normalsize Verified in Lean~4 / Mathlib~4 (v4.21.0)\\[2pt]
  \normalsize Repository: \href{https://github.com/JAgbanwa/miscellaneousmathproblems}%
    {\texttt{JAgbanwa/miscellaneousmathproblems}}
}
\date{\normalsize April 2026}

% ════════════════════════════════════════════════════════════════════════════
\begin{document}

\maketitle
\tableofcontents
\newpage

% ════════════════════════════════════════════════════════════════════════════
\section{Statement and Main Result}\label{sec:intro}

\subsection{The equation}

We study the Diophantine equation
\begin{equation}\label{eq:main}
  x^5 + y^4 = 3z^2, \qquad (x,y,z)\in\Z^3. \tag{$\star$}
\end{equation}

\subsection{Main theorem}

\begin{theorem}[Main result]\label{thm:main}
  The equation $x^5 + y^4 = 3z^2$ has \textbf{infinitely many} integer solutions.
  In particular, the primary parametric family
  \begin{equation}\label{eq:family1}
    (x, y, z) = \bigl(3a^2,\; 0,\; 9a^5\bigr), \qquad a \in \Z,
  \end{equation}
  provides infinitely many distinct integer solutions.
\end{theorem}

The proof is elementary: it requires only a \texttt{ring} computation to verify
the family~\eqref{eq:family1} and a strict-monotonicity argument to establish
that distinct non-negative~$a$ yield distinct solutions.
See Section~\ref{sec:proof} for the complete proof and Section~\ref{sec:lean}
for the Lean~4 formalisation.

% ════════════════════════════════════════════════════════════════════════════
\section{Derivation of the Parametric Families}\label{sec:derivation}

\subsection{Primary family via $y = 0$}\label{sec:primary}

Setting $y = 0$ reduces~\eqref{eq:main} to
\begin{equation}\label{eq:y0}
  x^5 = 3z^2.
\end{equation}
We seek integer solutions to \eqref{eq:y0}.  Write $x = 3a^2$ for $a \in \Z$:
\[
  x^5 = (3a^2)^5 = 3^5 a^{10} = 243 a^{10}.
\]
Then $3z^2 = 243a^{10}$ gives $z^2 = 81a^{10} = (9a^5)^2$, so $z = \pm 9a^5$. 
Choosing $z = 9a^5$ (the sign can always be absorbed into $a$), we obtain
\[
  (x, y, z) = (3a^2,\; 0,\; 9a^5), \qquad a \in \Z.
\]

\begin{lemma}[Primary family verification]\label{lem:primary}
  For every $a \in \Z$:
  \[
    (3a^2)^5 + 0^4 = 3 \cdot (9a^5)^2.
  \]
\end{lemma}
\begin{proof}
  $(3a^2)^5 = 3^5 a^{10} = 243 a^{10}$ and $3(9a^5)^2 = 3 \cdot 81 a^{10} = 243 a^{10}$. \qed
\end{proof}

The explicit table of small solutions from this family is:

\begin{center}
\begin{tabular}{rrrrl}
  \toprule
  $a$ & $x = 3a^2$ & $y$ & $z = 9a^5$ & Check\\
  \midrule
  $0$  & $0$   & $0$ & $0$       & $0=0$ \\
  $1$  & $3$   & $0$ & $9$       & $243 = 3\cdot81$ \\
  $-1$ & $3$   & $0$ & $-9$      & $243 = 3\cdot81$ \\
  $2$  & $12$  & $0$ & $288$     & $248832 = 3\cdot82944$ \\
  $-2$ & $12$  & $0$ & $-288$    & same \\
  $3$  & $27$  & $0$ & $2187$    & $14348907 = 3\cdot4782969$ \\
  \bottomrule
\end{tabular}
\end{center}

\subsection{Secondary family via weighted homogeneity}\label{sec:secondary}

The equation $x^5 + y^4 = 3z^2$ is invariant under the \emph{weighted scaling}
\[
  (x, y, z) \;\longmapsto\; (t^4\, x,\; t^5\, y,\; t^{10}\, z), \qquad t \in \Z.
\]
This follows from the \emph{weight assignment}
\[
  \mathrm{wt}(x) = 4,\quad \mathrm{wt}(y) = 5, \quad \mathrm{wt}(z) = 10,
\]
under which every monomial carries the same degree 20:
\[
  \mathrm{wt}(x^5) = 20, \quad \mathrm{wt}(y^4) = 20, \quad \mathrm{wt}(3z^2) = 20.
\]
Consequently, if $(x_0, y_0, z_0)$ is a solution, so is $(t^4 x_0, t^5 y_0, t^{10} z_0)$ for all $t \in \Z$.

\medskip
The triple $(2, 2, 4)$ is a seed solution: $2^5 + 2^4 = 32 + 16 = 48 = 3 \cdot 16 = 3 \cdot 4^2$. Applying the scaling gives:
\begin{equation}\label{eq:family2}
  (x, y, z) = \bigl(2t^4,\; 2t^5,\; 4t^{10}\bigr), \qquad t \in \Z.
\end{equation}

\begin{lemma}[Secondary family verification]\label{lem:secondary}
  For every $t \in \Z$:
  \[
    (2t^4)^5 + (2t^5)^4 = 3 \cdot (4t^{10})^2.
  \]
\end{lemma}
\begin{proof}
  $(2t^4)^5 + (2t^5)^4 = 32t^{20} + 16t^{20} = 48t^{20} = 3 \cdot 16 t^{20} = 3(4t^{10})^2$. \qed
\end{proof}

% ════════════════════════════════════════════════════════════════════════════
\section{Proof of Infinitely Many Solutions}\label{sec:proof}

\begin{proof}[Proof of Theorem~\ref{thm:main}]
  By Lemma~\ref{lem:primary}, every triple $(3a^2, 0, 9a^5)$ with $a \in \Z$ is a solution to~\eqref{eq:main}.

  It remains to show these are pairwise distinct.  Restricting to $a \in \N$
  (non-negative integers), the first component $x = 3a^2$ is a strictly
  increasing function of $a$:
  \[
    0 \leq a_1 < a_2 \;\Longrightarrow\; a_1^2 < a_2^2 \;\Longrightarrow\; 3a_1^2 < 3a_2^2.
  \]
  Thus for $a_1 \neq a_2$ in $\N$, the triples $(3a_1^2, 0, 9a_1^5)$ and $(3a_2^2, 0, 9a_2^5)$ have different $x$-coordinates and hence are distinct elements of $\Z^3$.

  Since the map $\N \to \Z^3$, $a \mapsto (3a^2, 0, 9a^5)$, is injective and its image lies in the solution set, the solution set is infinite.
\end{proof}

% ════════════════════════════════════════════════════════════════════════════
\section{Modular Analysis}\label{sec:modular}

\subsection{Mod 3}

By Fermat's little theorem, $x^5 \equiv x \pmod{3}$ and $y^4 \equiv y \pmod{3}$
(for $y \not\equiv 0$) or $y^4 \equiv 0$ (for $3 \mid y$).
Since $3z^2 \equiv 0 \pmod 3$, the equation mod 3 reduces to
\[
  x + y^4\bmod{3} \equiv 0 \pmod{3}.
\]
Two compatible cases arise:
\begin{enumerate}
  \item $3 \mid x$ and $3 \mid y$ (e.g., the primary family, since $x = 3a^2$), or
  \item $x \equiv 2 \pmod{3}$ and $3 \nmid y$ (e.g., the secondary family at $t=1$: $x = 2 \equiv 2$, $y = 2 \not\equiv 0$).
\end{enumerate}
No mod-3 obstruction to solutions.

\subsection{Parity}

$3z^2$ is odd iff $z$ is odd.  For $x^5 + y^4$: $y^4$ is even iff $y$ is even,
and $x^5 \equiv x \pmod{2}$.  Both parities are compatible (e.g., $x$ and $z$ both odd,
$y$ even in the primary family).  No parity obstruction.

\subsection{Summary}

No congruence obstruction prevents integer solutions.  This is consistent with the
existence of the two infinite parametric families above, and with the fact that the
equation has $\mathbb{F}_p$-points for every prime $p$ (verified computationally by
\texttt{brute\_force\_search.py}).

% ════════════════════════════════════════════════════════════════════════════
\section{Geometric Structure}\label{sec:geometry}

The equation $x^5 + y^4 = 3z^2$ defines a \emph{weighted projective hypersurface}
\[
  V \;\subset\; \mathbb{P}^2_{(4,5,10)}
\]
of weighted degree 20 in weighted projective space with weights $(4, 5, 10)$.  The affine
piece at $z=1$ is the surface $x^5 + y^4 = 3$.

Two explicit rational curves are embedded in $V$:
\begin{itemize}
  \item $C_1 \colon a \mapsto [3a^2 : 0 : 9a^5]$ (the primary family), which is a rational curve.
  \item $C_2 \colon t \mapsto [2t^4 : 2t^5 : 4t^{10}]$ (the secondary family), which is also rational.
\end{itemize}

Both curves are explicitly parametrised by $\mathbb{P}^1$, confirming that $V$ contains
rational curves and in particular has a Zariski-dense set of rational points
(infinitely many integer points on each curve).

This contrasts with the other Diophantine equations in this repository (for example,
$y^3 - x^4 = 0$ or $y^3 + xy + x^4 + 4 = 0$), which define algebraic curves of genus
$\geq 2$ and hence have only \emph{finitely} many rational points by Faltings' theorem.

% ════════════════════════════════════════════════════════════════════════════
\section{Lean~4 Formalisation}\label{sec:lean}

The file \texttt{InfinitelyManySolutions.lean} contains a complete Lean~4 / Mathlib~4
proof of Theorem~\ref{thm:main} with \textbf{zero \texttt{sorry} and zero additional axioms}.

\subsection{Summary of the formalisation}

\begin{center}
\begin{tabular}{lll}
  \toprule
  Lean name & Statement & Method \\
  \midrule
  \texttt{parametric\_solution} & $(3a^2)^5 + 0^4 = 3(9a^5)^2$ for all $a:\Z$ & \texttt{ring} \\
  \texttt{parametric\_solution2} & $(2t^4)^5 + (2t^5)^4 = 3(4t^{10})^2$ for all $t:\Z$ & \texttt{ring} \\
  \texttt{sol\_3\_0\_9} & $(3, 0, 9)$ is a solution & \texttt{norm\_num} \\
  \texttt{sol\_2\_2\_4} & $(2, 2, 4)$ is a solution & \texttt{norm\_num} \\
  \texttt{natMap\_mem} & Every $(3n^2, 0, 9n^5)$, $n:\N$, is a solution & \texttt{ring} \\
  \texttt{pow2\_strictMono} & $n \mapsto n^2$ is strictly monotone on $\N$ & \texttt{Nat.pow\_lt\_pow\_left} \\
  \texttt{natMap\_injective} & $n \mapsto (3n^2, 0, 9n^5)$ is injective & \texttt{mul\_left\_cancel\textsubscript{0}} \\
  \texttt{solutions\_infinite} & The solution set is infinite & \texttt{Set.infinite\_range\_of\_injective} \\
  \bottomrule
\end{tabular}
\end{center}

\subsection{Key Lean code excerpts}

\begin{lstlisting}[language=Lean4,caption={The equation and primary family verification}]
def isolution (x y z : Z) : Prop := x ^ 5 + y ^ 4 = 3 * z ^ 2

theorem parametric_solution (a : Z) :
    isolution (3 * a ^ 2) 0 (9 * a ^ 5) := by
  unfold isolution; ring
\end{lstlisting}

\begin{lstlisting}[language=Lean4,caption={Injectivity and infinitude}]
private theorem pow2_strictMono : StrictMono (fun n : N => n ^ 2) := by
  intro a b hab; exact Nat.pow_lt_pow_left hab (by norm_num)

theorem natMap_injective : Function.Injective natMap := by
  intro n1 n2 h
  simp only [natMap, Prod.mk.injEq] at h
  have h2 : (n1 : Z) ^ 2 = (n2 : Z) ^ 2 :=
    mul_left_cancel0 (by norm_num : (3 : Z) = 0) h.1
  have h3 : n1 ^ 2 = n2 ^ 2 := by exact_mod_cast h2
  exact pow2_strictMono.injective h3

theorem solutions_infinite : solutions.Infinite := by
  apply (Set.infinite_range_of_injective natMap_injective).mono
  rintro _ (n, rfl); exact natMap_mem n
\end{lstlisting}

\subsection{Build instructions}

The file is part of the Lake package in the repository root.  To verify:
\begin{lstlisting}[language=bash]
lake exe cache get
lake build
\end{lstlisting}

% ════════════════════════════════════════════════════════════════════════════
\section{Summary and Proof Status}\label{sec:summary}

\begin{center}
\begin{tabular}{lll}
  \toprule
  Component & Status & Method \\
  \midrule
  Seed solution $(3, 0, 9)$ & \checkmark\ Complete & Direct substitution \\
  Primary family $(3a^2, 0, 9a^5)$ verified & \checkmark\ Complete & \texttt{ring} identity \\
  Secondary family $(2t^4, 2t^5, 4t^{10})$ verified & \checkmark\ Complete & \texttt{ring} identity \\
  Infinitely many distinct solutions & \checkmark\ Complete & Strict monotonicity of $n^2$ on $\N$ \\
  Lean~4 formalisation (0 sorry) & \checkmark\ \textbf{Complete} & Mathlib tactics \\
  Full classification of all solutions & Open & Weighted projective geometry \\
  \bottomrule
\end{tabular}
\end{center}

\begin{theorem}[Final result]
  The Diophantine equation $x^5 + y^4 = 3z^2$ has \textbf{infinitely many} integer
  solutions.  The primary parametric family is
  \[
    (x, y, z) = (3a^2,\; 0,\; 9a^5), \qquad a \in \Z.
  \]
  A complete, unconditional Lean~4 proof with zero \texttt{sorry} and zero additional
  axioms is given in \texttt{InfinitelyManySolutions.lean}.
\end{theorem}

\end{document}
